\section{8.1 팀 프로젝트: 데이터 선정부터 발표까지}
\begin{frame}{왜 지금 프로젝트를 하는가}
  \begin{columns}[T]
    \begin{column}{0.55\textwidth}
      \begin{itemize}
        \item 트리 기반 모델(Week 7)까지 끝내면 정형 데이터의
          대표 도구를 \textbf{거의 다} 손에 쥔다 --- Block B 신경망
          없이도 실전 데이터셋 상당수를 다룰 수 있다.
        \item 지금이 ``각 장의 코드가 \textbf{실제 데이터} 앞에서도
          작동하는가''를 확인할 가장 좋은 시점.
      \end{itemize}
    \end{column}
    \begin{column}{0.43\textwidth}
      \kb{검증의 성격이 달라진다}
      \begin{itemize}
        \item Week 2--7 실습: 교재가 고른 데이터, 교재가 고른 모델.
        \item 이번엔 데이터 고르기, 모델 고르기, \textbf{나빠지면
          원인을 찾아 다시 고르기}까지가 우리 몫.
        \item ``코드가 동작한다''의 증거와 ``내가 풀 수 있다''의
          증거는 다르다.
      \end{itemize}
    \end{column}
  \end{columns}
\end{frame}
